\begin{proof}[Proof of \cref{obj:thm:tv-envelope-design}]
Let
\[
\kappa:=\frac{\sigma_0^2}{n},
\]
with the convention that this real quotient is total, so the quotient is zero when the natural population size is zero. Since \(\sigma_0^2\ge0\) and \(n\ge0\), this gives \(\kappa\ge0\). % lean: div_nonneg

We use this nonnegativity through the following elementary scaling rule for infima, applied twice below. Let \(S\subseteq\mathbb R\) be nonempty and bounded below, and let \(\kappa\ge0\). Then \(\kappa S=\{\kappa u:u\in S\}\) is again nonempty and bounded below, and
\[
\inf\{\kappa\,u:u\in S\}=\kappa\,\inf S .
\]
For \(\kappa>0\) this is the standard homogeneity of the infimum; for \(\kappa=0\) both sides equal \(0\). Each application below is preceded by a verification that the set being scaled is nonempty and bounded below. % lean: Real.sInf_smul_of_nonneg

\begin{enumerate}
\item
First, \(W_\beta(p)\) is nonempty. Indeed, the order hypothesis gives \(\beta\ge 1\), the theorem assumes \(k\ge \beta\), and \(p\in S_{k,q}\). Hence \cref{obj:lem:unbiased-weight-set-nonempty} applies and gives a weight vector \(w\in W_\beta(p)\). % lean: unbiased_weight_set_nonempty

\item
Fix now any \(w\in W_\beta(p)\). By linearity of the design expectation and by the defining mean-curve restriction \(\mathbb E_\pi[\bar Y_j]=m_P(p_j)\) in \(\mathcal P_\beta\) of \cref{obj:def:rollout-law-class},
\[
\mathbb E_\pi\!\left[\sum_{j=0}^k w_j\bar Y_j\right]
=
\sum_{j=0}^k w_j\,\mathbb E_\pi[\bar Y_j]
=
\sum_{j=0}^k w_j\,m_P(p_j).
\]
Since \(p\in S_{k,q}\), every node \(p_j\) lies in \([0,1]\), so the degree-\(\beta\) polynomial condition in \(\mathcal P_\beta\) gives
\[
m_P(p_j)=\sum_{\ell=0}^{\beta} a_{P,\ell}p_j^\ell .
\]
Substituting and interchanging the two finite sums,
\[
\sum_{j=0}^k w_jm_P(p_j)
=
\sum_{\ell=0}^{\beta}
a_{P,\ell}\sum_{j=0}^k w_jp_j^\ell .
\]
Splitting off the constant term and inserting the moment equations in \cref{obj:def:unbiased-weight-set},
\[
\sum_{j=0}^k w_jp_j^0=0,
\qquad
\sum_{j=0}^k w_jp_j^\ell=1
\quad(1\le \ell\le \beta),
\]
the last display equals
\[
a_{P,0}\sum_{j=0}^k w_jp_j^0
+
\sum_{\ell=1}^{\beta}a_{P,\ell}\sum_{j=0}^k w_jp_j^\ell
=
a_{P,0}\cdot 0+\sum_{\ell=1}^{\beta}a_{P,\ell}\cdot 1
=
\sum_{\ell=1}^{\beta}a_{P,\ell}.
\]
By \cref{obj:prop:rollout-polynomial-identity}, whose hypotheses are the static-rollout and \(\beta\)-polynomial members of \(\mathcal P_\beta\), this is \(m_P(1)-m_P(0)\), proving the unbiasedness identity. % lean: rollout_polynomial_identity

\item
For the variance statement, apply \cref{obj:lem:variance-envelope-sharpness} with \(X_j=\bar Y_j\). Its hypotheses are the nonnegative-scale assumption \(\sigma_0^2\ge0\) and the round-mean variance-envelope component of \(\mathcal P_\beta\). Therefore
\[
\operatorname{Var}_\pi\!\left(\sum_{j=0}^k w_j\bar Y_j\right)
\le
\frac{\sigma_0^2}{n}\left(\sum_{j=0}^k |w_j|\right)^2,
\]
and the same lemma supplies a positive semidefinite matrix \(\Gamma\) with
\[
\Gamma_{jj}\le \frac{\sigma_0^2}{n}
\quad\text{for all }j,
\qquad
\sum_{i=0}^k\sum_{j=0}^k w_i\Gamma_{ij}w_j
=
\frac{\sigma_0^2}{n}\left(\sum_{j=0}^k |w_j|\right)^2 .
\] % lean: variance_envelope_sharpness

\item
For the fixed schedule \(p\), the set whose infimum appears on the left-hand side of the first envelope identity is exactly the scalar multiple
\[
\left\{
v:\ \exists w\in W_\beta(p),\
v=\kappa\left(\sum_{j=0}^k |w_j|\right)^2
\right\}
=
\kappa\cdot\left\{u:\ \exists w\in W_\beta(p),\ u=\left(\sum_{j=0}^k |w_j|\right)^2\right\}.
\]
The unscaled set is nonempty by the first step and is bounded below by \(0\), since its elements are squares. Hence the usual nonnegative-scalar infimum rule applies to this set, and \cref{obj:def:amplification-criterion} identifies its infimum as \(A_\beta(p)\). With \(\kappa\ge0\),
\[
\inf\left\{
v:\ \exists w\in W_\beta(p),\
v=\frac{\sigma_0^2}{n}\left(\sum_{j=0}^k |w_j|\right)^2
\right\}
=
\frac{\sigma_0^2}{n}A_\beta(p).
\] % lean: Real.sInf_smul_of_nonneg

\item
The same scalar-infimum step applies after also varying the schedule. Namely,
\[
\left\{
v:\ \exists p'\in S_{k,q},\
v=\kappa\,A_\beta(p')
\right\}
=
\kappa\cdot\left\{u:\ \exists p'\in S_{k,q},\ u=A_\beta(p')
\right\}.
\]
This unscaled schedule set is nonempty because the fixed schedule \(p\) belongs to \(S_{k,q}\). It is bounded below by \(0\): for every schedule \(r\), \(A_\beta(r)\ge0\), because \(A_\beta(r)\) is the real infimum of squared total-variation norms, with the same convention covering the empty feasible-weight case. Applying the nonnegative-scalar infimum rule and then \cref{obj:def:amplification-criterion},
\[
\inf\left\{
v:\ \exists p'\in S_{k,q},\
v=\frac{\sigma_0^2}{n}A_\beta(p')
\right\}
=
\frac{\sigma_0^2}{n}M_{\beta,k,q}.
\] % lean: Real.sInf_smul_of_nonneg, amplification_nonneg

\item
Finally, suppose \(p^\star\in S_{k,q}\) and
\[
A_\beta(p^\star)=M_{\beta,k,q}.
\]
For any \(p'\in S_{k,q}\), the value \(A_\beta(p')\) belongs to the set
\[
\left\{u:\ \exists r\in S_{k,q},\ u=A_\beta(r)\right\}.
\]
The same uniform nonnegativity \(A_\beta(r)\ge0\) for all schedules \(r\) makes this set bounded below by \(0\). Therefore the defining infimum satisfies
\[
M_{\beta,k,q}\le A_\beta(p').
\]
Using \(A_\beta(p^\star)=M_{\beta,k,q}\), we get \(A_\beta(p^\star)\le A_\beta(p')\). Multiplication by the nonnegative scalar \(\kappa=\sigma_0^2/n\) preserves the inequality, so
\[
\frac{\sigma_0^2}{n}A_\beta(p^\star)
\le
\frac{\sigma_0^2}{n}A_\beta(p'),
\]
as required. % lean: minimaxAmplification_le_of_budgeted
\end{enumerate}
\end{proof}

\begin{proof}[Proof of \cref{obj:prop:equal-spacing-benchmark}]
Let \(p_j=qj/\beta\) for \(j=0,\ldots,\beta\). Since \(\beta\ge1\), \(q>0\), and \(q\le1\), each coordinate is nonnegative, and \(j\le\beta\) gives
\[
0\le p_j=\frac{qj}{\beta}\le q\le1.
\]
Thus \(p_j\in[0,1]\) for every \(j=0,\ldots,\beta\). % lean: hrange

For the amplification bound, use the Lagrange endpoint weights on the equal grid. The grid is injective, since \(p_j=p_m\) forces \(qj/\beta=qm/\beta\) and hence \(j=m\) because \(q>0\) and \(\beta>0\); % lean: equalSchedule_injective
so the Lagrange basis polynomials
\[
L_j(t)=\prod_{\substack{0\le m\le\beta\\ m\ne j}}\frac{t-p_m}{p_j-p_m},
\qquad j=0,\ldots,\beta,
\]
are well defined, and we set
\[
w_j=L_j(1)-L_j(0),
\qquad j=0,\ldots,\beta .
\]
By Lagrange interpolation on the \(\beta+1\) distinct nodes, these weights reproduce \(r(1)-r(0)\) for every polynomial \(r\) of degree at most \(\beta\); applied to the monomials \(1,u,\ldots,u^\beta\), this is exactly \(w\in W_\beta(p)\) as in \cref{obj:def:unbiased-weight-set}. % lean: lagrange_endpoint_weights_unbiased
Since \(A_\beta(p)\) is the infimum in \cref{obj:def:amplification-criterion} over the nonempty set \(W_\beta(p)\), this feasible \(w\) gives
\[
A_\beta(p)\le \left(\sum_{j=0}^{\beta}|w_j|\right)^2.
\] % lean: hA_le

Set \(y=\beta/q\). Since \(q\le1\le\beta\), we have \(y\ge1\), and hence \(y^\beta\ge1\). % lean: hy_ge_one

\emph{The value at \(1\).}
Fix \(j\). Each factor of \(L_j(1)\) is controlled separately: from \(0\le p_m\le q\le 1\) we get \(|1-p_m|\le1\), while
\[
p_j-p_m=\frac{q}{\beta}(j-m),
\qquad\text{so}\qquad
|p_j-p_m|=\frac{q}{\beta}\,|j-m| .
\]
There are exactly \(\beta\) indices \(m\ne j\) in \(\{0,\ldots,\beta\}\), so multiplying the \(\beta\) factor bounds gives
\[
|L_j(1)|
=
\prod_{m\ne j}\frac{|1-p_m|}{|p_j-p_m|}
\le
\prod_{m\ne j}\frac{\beta/q}{|j-m|}
=
\frac{y^\beta}{\prod_{m\ne j}|j-m|}.
\]
The remaining product is a factorial pair: splitting the indices below and above \(j\),
\[
\prod_{\substack{0\le m\le\beta\\ m\ne j}}|j-m|
=
\left(\prod_{m=0}^{j-1}(j-m)\right)\left(\prod_{m=j+1}^{\beta}(m-j)\right)
=
j!\,(\beta-j)! .
\]
Hence
\[
|L_j(1)|
\le
\frac{y^\beta}{j!(\beta-j)!}.
\] % lean: equalSchedule_lagrange_basis_eval_one_le
Summing over \(j\) requires the binomial identity
\[
\sum_{j=0}^{\beta}\frac{1}{j!(\beta-j)!}
=
\frac{1}{\beta!}\sum_{j=0}^{\beta}\frac{\beta!}{j!(\beta-j)!}
=
\frac{2^\beta}{\beta!}
\le
2,
\]
where the first equality rescales each term by \(\beta!\), the second sums the binomial coefficients \(\beta!/(j!(\beta-j)!)\) over \(j=0,\ldots,\beta\), and the last inequality is \(2^\beta\le2\,\beta!\), valid for every \(\beta\ge1\) by induction (the base case \(\beta=1\) reads \(2\le2\), and the step multiplies the left side by \(2\) and the right side by \(\beta+1\ge2\)). % lean: factorial_reciprocal_sum_le_two
Therefore
\[
\sum_{j=0}^{\beta}|L_j(1)|
\le
y^\beta\sum_{j=0}^{\beta}\frac{1}{j!(\beta-j)!}
\le
2y^\beta .
\]

\emph{The value at \(0\).}
The node \(p_0\) is exactly \(0\), so \(L_j(0)=L_j(p_0)\) is the Lagrange basis polynomial evaluated at a node: it equals \(1\) for \(j=0\) and \(0\) for \(j\ne0\). Hence
\[
\sum_{j=0}^{\beta}|L_j(0)|=1.
\] % lean: hsum_eval0

\emph{Conclusion.}
By the triangle inequality applied termwise to \(w_j=L_j(1)-L_j(0)\),
\[
\sum_{j=0}^{\beta}|w_j|
\le
\sum_{j=0}^{\beta}|L_j(1)|+\sum_{j=0}^{\beta}|L_j(0)|
\le
2y^\beta+1
\le
3y^\beta,
\]
where the last inequality uses \(y^\beta\ge1\). % lean: hsumw
Both sides are nonnegative, so squaring and substituting \(y=\beta/q\) yields
\[
A_\beta\!\left(p^{\mathrm{eq}}(\beta,q)\right)
\le
\left(\sum_{j=0}^{\beta}|w_j|\right)^2
\le
(3y^\beta)^2
=
9\left(\frac{\beta}{q}\right)^{2\beta}.
\] % lean: equal_spacing_benchmark
\end{proof}

\begin{proof}[Proof of \cref{obj:prop:no-extrapolation-boundary}]
At the full budget \(q=1\) and with \(k=\beta\), the equal-spacing schedule has entries
\[
p^{\mathrm{eq}}_j(\beta,1)=\frac{j}{\beta},
\qquad j=0,\ldots,\beta,
\qquad\text{so in particular}\qquad
p^{\mathrm{eq}}_0(\beta,1)=0,
\quad
p^{\mathrm{eq}}_\beta(\beta,1)=1 .
\]
Write \(p_j=p^{\mathrm{eq}}_j(\beta,1)\) and
\[
w_0=-1,\qquad w_\beta=1,\qquad w_j=0\quad(1\le j\le\beta-1).
\]
Because \(\beta\ge1\), the first and last grid indices are distinct, i.e. \(\beta\ne0\), so the two prescriptions \(w_0=-1\) and \(w_\beta=1\) do not conflict and the vector \(w\) is well defined. % lean: hne

\emph{The moment equations.}
Only the indices \(j=0\) and \(j=\beta\) carry nonzero weight. For the zeroth moment, \(p_j^0=1\) for every \(j\) (including \(p_0=0\)), so
\[
\sum_{j=0}^{\beta}w_j p_j^0
=
w_0+w_\beta
=
(-1)+1
=
0.
\] % lean: hw
Now fix \(1\le\ell\le\beta\). Since \(p_0=0\) gives \(p_0^\ell=0^\ell=0\) and \(p_\beta=1\) gives \(p_\beta^\ell=1\),
\[
\sum_{j=0}^{\beta}w_j p_j^\ell
=
w_0\,p_0^\ell+w_\beta\,p_\beta^\ell
=
-0^\ell+1^\ell
=
1,
\]
with all interior terms vanishing because their weights are zero. These are exactly the moment equations of \cref{obj:def:unbiased-weight-set}, so \(w\in W_\beta(p^{\mathrm{eq}}(\beta,1))\). % lean: no_extrapolation_boundary

\emph{The amplification bound.}
By \cref{obj:def:amplification-criterion}, \(A_\beta(p^{\mathrm{eq}}(\beta,1))\) is the infimum of \(\left(\sum_j|w_j'|\right)^2\) over \(w'\in W_\beta(p^{\mathrm{eq}}(\beta,1))\), a set of nonnegative numbers, so the feasible \(w\) just exhibited gives
\[
A_\beta\!\left(p^{\mathrm{eq}}(\beta,1)\right)
\le
\left(\sum_{j=0}^{\beta}|w_j|\right)^2.
\] % lean: csInf_le
The absolute values contribute only at \(j=0\) and \(j=\beta\), so
\[
\sum_{j=0}^{\beta}|w_j|=|-1|+|1|=1+1=2.
\] % lean: hsum_abs
Therefore
\[
A_\beta\!\left(p^{\mathrm{eq}}(\beta,1)\right)\le 2^2=4.
\] % lean: no_extrapolation_boundary
\end{proof}

\begin{proof}[Proof of \cref{obj:thm:chebyshev-minimax}]
\begin{enumerate}
\item \emph{The constants.}
By \cref{obj:lem:continuous-chebyshev-endpoint-bound}, applied with the cap \(q_{\max}\), choose constants
\(C_{\mathrm{ep}}(q_{\max})>0\) and \(c_{\mathrm{ep}}(q_{\max})>0\) such that, writing
\(x_q=2/q-1\) and
\[
\lambda(q)=x_q+\sqrt{x_q^2-1},
\]
the following two bounds hold for every \(\beta\ge1\) and every \(0<q\le q_{\max}\): first, every real polynomial \(R\) with \(\deg(R)\le\beta\) and \(\sup_{x\in[-1,1]}|R(x)|\le1\) satisfies
\[
|R(x_q)-R(-1)|\le C_{\mathrm{ep}}(q_{\max})\,\lambda(q)^\beta ;
\]
and second,
\[
T_\beta(x_q)-1\ge c_{\mathrm{ep}}(q_{\max})\,\lambda(q)^\beta .
\]
These constants are chosen before the oversampling ratio is fixed, so set
\[
C_-(q_{\max})=c_{\mathrm{ep}}(q_{\max})^2>0.
\]
Now fix \(c>1\). The Ehlich--Zeller Chebyshev--Lobatto mesh inequality used as the premise of \cref{obj:lem:oversampled-chebyshev-lobatto-norming} is available in the following form: for every pair of nonnegative integers \(\beta,k\), every real polynomial \(R\) with \(\deg(R)\le\beta\) and \(\beta<k\), and every \(x\in[-1,1]\),
\[
 |R(x)|
 \le
 \frac{1}{\cos(\pi\beta/(2k))}
 \max_{0\le j\le k}\left|R\!\left(-\cos\!\left(\frac{\pi j}{k}\right)\right)\right|.
\]
% lean: ehlichZellerMesh
Applying \cref{obj:lem:oversampled-chebyshev-lobatto-norming} with this mesh inequality, choose
\(K(c)>0\) with the property that every real polynomial \(R\) of degree at most \(\beta\) obeys
\(\sup_{x\in[-1,1]}|R(x)|\le K(c)\max_{0\le j\le k}\left|R(-\cos(\pi j/k))\right|\) whenever \(\beta\ge1\) and \(k\ge c\beta\), and set
\[
C_+(c,q_{\max})=\bigl(K(c)C_{\mathrm{ep}}(q_{\max})\bigr)^2>0 .
\]
% lean: chebyshev_minimax

\item \emph{Standing consequences of the hypotheses.}
Fix \(\beta,k,q\) satisfying the hypotheses. The low-budget-regime hypothesis gives \(q>0\), and the low-budget cap of \cref{obj:ass:low-budget-cap} gives
\(q\le q_{\max}<1\); hence \(q\le 1\), \(q<1\), and \(0<q\le 1\). Since \(k\ge c\beta\), \(c>1\), and
\(\beta\ge1\), we have \(\beta\le c\beta\le k\), and in particular \(k\ge1\). Therefore
\cref{obj:lem:chebyshev-schedule-admissible} gives
\[
p^{\mathrm{Ch}}(k,q)\in S_{k,q}.
\]
% lean: chebyshev_schedule_admissible
The exterior parameter is the exponential base in the statement:
\[
\lambda(q)=\frac{(1+\sqrt{1-q})^2}{q}.
\]
Indeed, \(x_q^2-1=(2/q-1)^2-1=4/q^2-4/q=4(1-q)/q^2\), which is nonnegative because \(q<1\), so \(\sqrt{x_q^2-1}=2\sqrt{1-q}/q\); adding \(x_q=(2-q)/q\) gives \(\lambda(q)=\bigl(2-q+2\sqrt{1-q}\bigr)/q\), and expanding \((1+\sqrt{1-q})^2=1+2\sqrt{1-q}+(1-q)=2-q+2\sqrt{1-q}\) identifies the two expressions.
% lean: chebyshev_lambda_eq_rho_div_q
Finally, for any schedule \(p'\in S_{k,q}\) we record the dual quantity of \cref{obj:lem:amplification-dual-norm},
\[
D(p')=\sup\left\{|r(1)-r(0)|:
r\in\mathbb R[x],\ \deg(r)\le\beta,\ \max_{0\le j\le k}|r(p'_j)|\le1\right\},
\]
and note that \(D(p')\ge0\) because the zero polynomial is feasible.

\item \emph{Lower bound at an arbitrary \(p\in S_{k,q}\).}
Since \(\beta\ge1\), \(\beta\le k\), and \(p\in S_{k,q}\), \cref{obj:lem:unbiased-weight-set-nonempty}
makes \(W_\beta(p)\) nonempty; and the strict ordering \(p_0<\cdots<p_k\) required by
\cref{obj:def:budgeted-schedule-class} makes the nodes distinct. Thus
\cref{obj:lem:amplification-dual-norm} applies and gives
\[
A_\beta(p)=D(p)^2 .
\]
% lean: amplification_dual_norm

Consider the test polynomial
\[
r(u)=T_\beta\!\left(\frac{2u}{q}-1\right).
\]
Its degree is at most \(\beta\), being the composition of the degree-\(\beta\) polynomial \(T_\beta\) with an affine map. % lean: chebyshev_affine_natDegree_le
For every rollout node, \(p\in S_{k,q}\) gives \(p_j\ge p_0=0\) and, by monotonicity, \(p_j\le p_k=q\); hence
\[
0\le p_j\le q,
\qquad\text{so}\qquad
\frac{2p_j}{q}-1\in[-1,1],
\]
and since \(|T_\beta|\le1\) on \([-1,1]\) we get \(|r(p_j)|\le1\). % lean: budgetedSchedule_le_endpoint
Therefore \(r\) is feasible in the supremum defining \(D(p)\). Its endpoint values are
\[
r(1)=T_\beta\!\left(\frac2q-1\right)=T_\beta(x_q),
\qquad
r(0)=T_\beta(-1).
\]
The endpoint lower bound from \cref{obj:lem:continuous-chebyshev-endpoint-bound} gives
\[
T_\beta(x_q)-1\ge c_{\mathrm{ep}}(q_{\max})\lambda(q)^\beta .
\]
Since \(x_q>1\) we have \(T_\beta(x_q)\ge1>0\), so \(|T_\beta(x_q)|=T_\beta(x_q)\); and \(|T_\beta(-1)|\le1\). The reverse triangle inequality therefore yields
\[
T_\beta(x_q)-1
\le
|T_\beta(x_q)|-|T_\beta(-1)|
\le |T_\beta(x_q)-T_\beta(-1)|
=|r(1)-r(0)|
\le D(p).
\]
Hence \(D(p)\ge c_{\mathrm{ep}}(q_{\max})\lambda(q)^\beta\ge0\), and squaring both nonnegative sides,
\[
A_\beta(p)=D(p)^2\ge c_{\mathrm{ep}}(q_{\max})^2\lambda(q)^{2\beta}
=
C_-(q_{\max})
\left(\frac{(1+\sqrt{1-q})^2}{q}\right)^{2\beta}.
\]
% lean: chebyshev_amplification_lower

\item \emph{Upper bound at the Chebyshev--Lobatto schedule.}
The schedule \(p^{\mathrm{Ch}}(k,q)\) is admissible by \cref{obj:lem:chebyshev-schedule-admissible}. Together with \(\beta\ge1\) and \(\beta\le k\), \cref{obj:lem:unbiased-weight-set-nonempty} makes \(W_\beta(p^{\mathrm{Ch}}(k,q))\) nonempty; the strict ordering in \(p^{\mathrm{Ch}}(k,q)\in S_{k,q}\) gives distinct nodes. Thus \cref{obj:lem:amplification-dual-norm} gives
\(A_\beta(p^{\mathrm{Ch}}(k,q))=D(p^{\mathrm{Ch}}(k,q))^2\), and it suffices to bound the dual supremum. Let \(r\) be
any polynomial with \(\deg(r)\le\beta\) and
\(\max_{0\le j\le k} |r(p^{\mathrm{Ch}}_j(k,q))|\le1\), and put
\[
R(x)=r\!\left(\frac q2(x+1)\right).
\]
Then \(\deg(R)\le\beta\), again because the inner map is affine. % lean: lobatto_affine_comp_natDegree_le
At the Lobatto abscissae the inner map reproduces the schedule,
\[
\frac q2\left(-\cos\frac{\pi j}{k}+1\right)
=
\frac q2\left(1-\cos\frac{\pi j}{k}\right)
=
p^{\mathrm{Ch}}_j(k,q),
\qquad j=0,\ldots,k,
\]
so \(R(-\cos(\pi j/k))=r(p^{\mathrm{Ch}}_j(k,q))\) and therefore
\[
\max_{0\le j\le k}\left|R\!\left(-\cos\frac{\pi j}{k}\right)\right|\le1 .
\]
The oversampled norming inequality supplied by \cref{obj:lem:oversampled-chebyshev-lobatto-norming} with the displayed Ehlich--Zeller mesh input applies because \(\beta\ge1\), \(k\ge c\beta\), and \(\deg(R)\le\beta\). It gives
\[
|R(x)|\le K(c)\cdot 1=K(c)\qquad (x\in[-1,1]).
\]
Thus \(S=K(c)^{-1}R\) satisfies \(\deg(S)\le\beta\) and \(|S(x)|\le1\) on \([-1,1]\). The inner affine map sends the two relevant points to the two endpoints,
\[
\frac q2\left(x_q+1\right)=\frac q2\cdot\frac2q=1,
\qquad
\frac q2\left(-1+1\right)=0,
\]
so \(R(x_q)=r(1)\) and \(R(-1)=r(0)\), and hence \(S(x_q)=K(c)^{-1}r(1)\) and \(S(-1)=K(c)^{-1}r(0)\). Applying the endpoint upper bound from
\cref{obj:lem:continuous-chebyshev-endpoint-bound} to \(S\) gives
\[
K(c)^{-1}|r(1)-r(0)|
=
|S(x_q)-S(-1)|
\le C_{\mathrm{ep}}(q_{\max})\lambda(q)^\beta,
\]
that is, \(|r(1)-r(0)|\le K(c)C_{\mathrm{ep}}(q_{\max})\lambda(q)^\beta\). Taking the supremum over all such \(r\) gives
\(0\le D(p^{\mathrm{Ch}}(k,q))\le K(c)C_{\mathrm{ep}}(q_{\max})\lambda(q)^\beta\), and squaring both nonnegative sides,
\[
A_\beta\!\left(p^{\mathrm{Ch}}(k,q)\right)
\le
\bigl(K(c)C_{\mathrm{ep}}(q_{\max})\bigr)^2\lambda(q)^{2\beta}
=
C_+(c,q_{\max})
\left(\frac{(1+\sqrt{1-q})^2}{q}\right)^{2\beta}.
\]
% lean: chebyshev_amplification_upper

\item \emph{The two-sided minimax bound.}
The first two displayed inequalities of the statement are exactly the bounds just proved. For the minimax
lower bound, the universal lower estimate obtained from \cref{obj:lem:unbiased-weight-set-nonempty}, \cref{obj:lem:amplification-dual-norm}, and \cref{obj:lem:continuous-chebyshev-endpoint-bound} holds for every \(p\in S_{k,q}\), and the schedule class is
nonempty because \(p^{\mathrm{Ch}}(k,q)\in S_{k,q}\); so the common bound is a lower bound for the nonempty set \(\{A_\beta(p):p\in S_{k,q}\}\), and passing to its infimum gives
\[
M_{\beta,k,q}
\ge
C_-(q_{\max})
\left(\frac{(1+\sqrt{1-q})^2}{q}\right)^{2\beta}.
\]
% lean: minimaxAmplification_lower_of_forall
For the minimax upper bound, that set is bounded below by \(0\) and contains \(A_\beta(p^{\mathrm{Ch}}(k,q))\), so its infimum is at most that value; chaining with the Chebyshev upper bound obtained from \cref{obj:lem:chebyshev-schedule-admissible}, \cref{obj:lem:oversampled-chebyshev-lobatto-norming}, and \cref{obj:lem:continuous-chebyshev-endpoint-bound},
\[
M_{\beta,k,q}
\le
A_\beta\!\left(p^{\mathrm{Ch}}(k,q)\right)
\le
C_+(c,q_{\max})
\left(\frac{(1+\sqrt{1-q})^2}{q}\right)^{2\beta}.
\]
This proves the claimed two-sided minimax estimate.
% lean: minimaxAmplification_le_of_budgeted
\end{enumerate}
\end{proof}

\begin{proof}[Proof of \cref{obj:lem:exact-chebyshev-rate-feasible}]
By \cref{obj:thm:chebyshev-minimax}, applied with the fixed cap \(q_{\max}\) and the oversampling constant \(c>1\), there is a constant \(C_+(c,q_{\max})>0\) such that, whenever \(\beta\ge1\), \(k\ge c\beta\), \(q>0\), and the low-budget cap holds,
\[
A_\beta\!\left(p^{\mathrm{Ch}}(k,q)\right)
\le
C_+(c,q_{\max})
\left(\frac{(1+\sqrt{1-q})^2}{q}\right)^{2\beta}.
\]
% lean: chebyshev_minimax
This is the constant used in the statement; it depends only on \(c\) and \(q_{\max}\).

Now fix \(n,\beta,k,\Omega,\pi,q,\sigma_0^2\) satisfying the hypotheses of the lemma, and write \(p^{\mathrm{Ch}}=p^{\mathrm{Ch}}(k,q)\).

\emph{The hypotheses of the two inputs hold.}
Since \(k=\lceil c\beta\rceil\) is the least integer at or above \(c\beta\), and since \(c>1\) and \(\beta\ge1\) give \(\beta\le c\beta\), we have
\[
\beta\le c\beta\le k .
\] % lean: Nat.le_ceil
In particular \(k\ge\beta\ge1\), so \(k\ge1\). The low-budget cap of \cref{obj:ass:low-budget-cap} gives \(q\le q_{\max}<1\), which together with \(q>0\) yields
\[
0<q\le1 .
\] % lean: LowBudgetCap
Therefore \cref{obj:lem:chebyshev-schedule-admissible} applies and gives
\[
p^{\mathrm{Ch}}\in S_{k,q}.
\]
% lean: chebyshev_schedule_admissible
Finally, write \(s=\sigma_0^2/n\), interpreting the quotient as the total real division used in the risk definitions; in particular the quotient is defined also when \(n=0\), where it equals \(0\). Since \(\sigma_0^2\ge0\) and the real cast of a natural number is nonnegative, this gives
\[
s=\frac{\sigma_0^2}{n}\ge0
\]
uniformly over the whole range of population sizes, with no case distinction. % lean: div_nonneg

\emph{Envelope comparison at the Chebyshev schedule.}
For any weight vector \(w\), the exact quadratic form is by definition the design variance of the corresponding linear rollout estimator,
\[
w^\prime\Gamma_P\!\left(p^{\mathrm{Ch}}\right)w
=
\operatorname{Var}_\pi\!\left(\sum_{j=0}^k w_j\bar Y_j\right),
\]
because \(\Gamma_P(p)\) is the design covariance matrix of \((\bar Y_0,\ldots,\bar Y_k)\). Applying the fixed-schedule part of \cref{obj:lem:exact-risk-envelope-upper} to the admissible schedule \(p^{\mathrm{Ch}}\), whose hypotheses \(\beta\ge1\), \(k\ge\beta\), \(0<q\le1\), \(\sigma_0^2\ge0\) and \(p^{\mathrm{Ch}}\in S_{k,q}\) were just verified, gives
\[
\inf_{w\in W_\beta(p^{\mathrm{Ch}})}
\sup_{P\in\mathcal P_\beta}
\operatorname{Var}_\pi\!\left(\sum_{j=0}^k w_j\bar Y_j\right)
\le
\frac{\sigma_0^2}{n}
A_\beta\!\left(p^{\mathrm{Ch}}\right).
\]
% lean: exact_risk_envelope_upper

\emph{Conclusion.}
Since \(s=\sigma_0^2/n\ge0\), multiplying the Chebyshev upper bound for \(A_\beta(p^{\mathrm{Ch}})\) by \(\sigma_0^2/n\) preserves the inequality, so chaining the last two displays gives
\[
\inf_{w\in W_\beta(p^{\mathrm{Ch}})}
\sup_{P\in\mathcal P_\beta}
\operatorname{Var}_\pi\!\left(\sum_{j=0}^k w_j\bar Y_j\right)
\le
\frac{\sigma_0^2}{n}
A_\beta\!\left(p^{\mathrm{Ch}}\right)
\le
\frac{\sigma_0^2}{n}\,
C_+(c,q_{\max})
\left(\frac{(1+\sqrt{1-q})^2}{q}\right)^{2\beta},
\]
which is the asserted bound.
% lean: mul_le_mul_of_nonneg_left
\end{proof}
